% ============================================================
% sec/01introduccion.tex — VERSIÓN INTEGRADA, TRES EJES
%
% Eje A: Duarte · Eje B: Chaves · Eje C: Hernández
% Cifras procedentes de la auditoría del corpus.
%
% OBLIGATORIO: contexto y relevancia, pregunta de investigación
% explícita, hipótesis de trabajo, objetivos con contribución
% esperada, y estructura del documento.
% ============================================================
\section{Introducción}
\label{sec:introduccion}

\subsection{Contexto y relevancia}
\label{subsec:contexto}

Estimar la probabilidad de que un jugador venza desde la posición que ocupa
actualmente sobre la mesa es una competencia recurrente de los sistemas que
juegan. Una función de evaluación de estados convierte una búsqueda
intratable en una tratable: permite a un agente detener una simulación antes
de tiempo y aun así devolver un juicio con sentido. En los juegos
competitivos esa misma estimación tiene un segundo público, pues analistas y
jugadores la leen como una afirmación sobre qué recursos deciden realmente
una partida. Ambos usos dependen del mismo objeto: una función que asigna a
una posición observable una probabilidad de victoria.

Esa cantidad posee una propiedad que conviene fijar desde el comienzo. Brill
\textit{et al.} señalan que la probabilidad de victoria no es una estadística
contable ni una magnitud observable, sino que queda definida por un modelo
estimado a partir de datos \cite{brill_winprob_2026}. El desenlace observado
de una partida es binario, mientras que la magnitud que se desea estimar es
una probabilidad que nunca se observa de forma directa.

\razon{Esta declaración debe aparecer en los primeros párrafos. Si no se
establece que la probabilidad verdadera es inobservable, el lector puede
suponer que existe una etiqueta contra la cual verificar cada predicción, y
toda la discusión posterior sobre calibración pierde sentido.}

\subsubsection{Dos propiedades que dificultan la tarea en juegos de cartas}

La primera es la información imperfecta. Un jugador nunca ve la mano del
oponente ni las cartas de premio, de modo que el evaluador opera sobre una
vista parcial y no sobre el estado verdadero del mundo.

La segunda es la no estacionariedad. El repertorio de cartas legales, los
mazos que los jugadores llevan y la población de adversarios cambian con el
tiempo, y un modelo ajustado a un conjunto de estrategias dominante no tiene
por qué sobrevivir al siguiente. Bertram \textit{et al.} lo muestran de forma
concreta para \emph{Magic: The Gathering}: una representación de carta
basada en codificación por identidad alcanza un 67.80\,\% de acierto sobre
el conjunto de cartas con el que fue entrenada y queda simplemente indefinida
sobre cartas que nunca vio, mientras que una representación construida a
partir de atributos conserva un 68.00\,\% dentro de la distribución y aún
alcanza un 42.87\,\% sobre conjuntos no vistos
\cite{bertram_cardrepresentations_2024}.

\subsubsection{La tarea en el dominio de los juegos competitivos}

En deportes electrónicos de tipo arena multijugador, la literatura trata la
predicción de resultado como una aplicación consolidada del aprendizaje
automático, y emplea tanto variables previas a la partida como información
observada durante ella \cite{kamal_moba_2025}. Tian \textit{et al.} realizan
predicción en tiempo real a partir de secuencias de estados y muestran que la
importancia de la alineación, el oro, las bajas y las torres varía con el
progreso del encuentro \cite{tian_hero_2022}.

La utilidad de esta estimación excede el interés metodológico. Hodge
\textit{et al.} observan que, a diferencia de los deportes tradicionales,
numerosos juegos competitivos carecen de un marcador legible, y que su
complejidad dificulta el seguimiento por parte de audiencias no expertas; una
estadística única que resuma el estado de la partida amplía el atractivo del
juego y lo vuelve más accesible \cite{hodge_winprediction_2021}. Xenopoulos
\textit{et al.} desarrollan la misma línea al estimar la probabilidad de
victoria por ronda y analizar cómo varían los efectos de las características
del estado entre poblaciones de jugadores con distinto nivel de habilidad
\cite{xenopoulos_winprob_2022}.

En juegos de cartas la idea aparece con objetivos distintos. Miernik y
Kowalski evolucionan funciones de evaluación para \emph{Legends of Code and
Magic}, cuya salida es un puntaje heurístico optimizado de forma indirecta
por el rendimiento del agente y no una probabilidad calibrada de victoria
\cite{miernik_evolving_2022}. Más cerca de la formulación adoptada aquí, el
AAIA'17 Data Mining Challenge solicitó predecir el ganador de una partida de
\emph{Hearthstone} a partir de un único estado de juego, y reportó un
resultado oficial de referencia de 0.7846 de área bajo la curva ROC frente a
una mejor solución de 0.8019 \cite{janusz_hearthstone_2017}. Para
\emph{StarCraft~II}, el descriptor SC2EGSet acompaña su corpus de
repeticiones con experimentos de predicción de victoria sobre indicadores
económicos agregados, y alcanza un acierto de 0.9071 cuando dispone de las
características de ambos jugadores \cite{bialecki_sc2egset_2023}.

El \emph{Pok\'emon Trading Card Game} carece de un estudio publicado
comparable, pese a que su estructura de recursos, cartas de premio,
asignación de energía, banca, líneas de evolución, difiere de la de ambos
juegos de maneras que deberían cambiar qué características transportan
información.

\razon{Declarar que la tarea no es inédita fortalece el trabajo en lugar de
debilitarlo. Un planteamiento que reclamara originalidad absoluta sería
contradicho por las referencias del propio documento, y la contribución real
se sitúa en el traslado al dominio y en el tratamiento metodológico.}

\subsubsection{Generalización ante contenido y poblaciones cambiantes}

Chitayat \textit{et al.} muestran que representaciones ligadas a parámetros
de diseño pueden conservar su rendimiento al atravesar actualizaciones
posteriores del juego, mientras que las basadas únicamente en identidad
resultan más frágiles \cite{chitayat_meta_2023}. En el videojuego competitivo
de la franquicia Pok\'emon, Angliss \textit{et al.} separan de forma
explícita el rendimiento sobre composiciones vistas de la generalización a
composiciones no vistas, y muestran que entrenar con mayor diversidad mejora
esa segunda capacidad \cite{angliss_vgcbench_2026}.

Estos antecedentes motivan que una evaluación sobre el \emph{Pok\'emon TCG}
no se limite a una partición aleatoria de partidas, sino que examine además
cambios temporales y composiciones no observadas durante el entrenamiento.

\pc{Dos o tres oraciones con datos concretos del dominio del juego de cartas:
tamaño de la comunidad competitiva, cantidad de cartas disponibles,
complejidad combinatoria de la construcción de mazos.}

El estado de una partida involucra premios restantes, puntos de vida,
energías acumuladas, cartas en mano, evoluciones y estados alterados,
distribuidos entre dos jugadores con información parcial sobre el mazo y la
mano del oponente. Un observador no experto no puede determinar de un vistazo
cuál de los dos se encuentra en ventaja.

\subsection{Planteamiento del problema}
\label{subsec:problema}

Este trabajo aborda una limitación específica: para el \emph{Pok\'emon TCG}
no existe evidencia sobre cuánta información contiene un único estado
observable acerca del desenlace, sobre cómo se acumula esa información a lo
largo de la partida, ni sobre cuánta se transfiere a partidas, mazos y
periodos temporales no vistos durante el entrenamiento.

La literatura que respondería esas preguntas para juegos adyacentes responde
solo una parte. Janusz \textit{et al.} reportan discriminación pero no
calibración \cite{janusz_hearthstone_2017}; Bia\l ecki \textit{et al.}
reportan acierto y describen sus propios experimentos como más ilustrativos
que investigativos \cite{bialecki_sc2egset_2023}; y ninguno reporta
rendimiento bajo un cambio de distribución declarado.

\subsubsection{La estructura del corpus y su consecuencia}

Esa brecha se ha vuelto abordable recientemente. El PTCG AI Battle Challenge
publica a diario repeticiones en formato JSON de partidas completas entre
agentes que compiten entre sí \cite{ptcg_dataset_2026}. El corpus declara 67
entregas diarias, del 2026-06-16 al 2026-08-21, con 329\,336 episodios y
1.29~TB sin comprimir.

Nuestra auditoría cubre una muestra temporal de diez de esos días, el
14.9\,\% de los días y el 15.8\,\% de los episodios declarados. Dentro de esa
muestra se analizaron 52\,085 episodios a partir de 182.6~GiB de JSON con
cero errores de análisis sintáctico, y se obtuvieron 7\,038\,378
instantáneas de estado utilizables: aproximadamente 135 filas correlacionadas
por cada partida independiente.

Datos con esta forma premian el protocolo equivocado. Si esas instantáneas se
particionaran al azar, estados de una misma partida aparecerían
simultáneamente en entrenamiento y en prueba. Lo medimos directamente: bajo
una partición aleatoria 70/30 a nivel de instantánea, la fracción esperada de
episodios presentes en ambas particiones es 0.999945.

Brill \textit{et al.} demuestran que, cuando múltiples observaciones
comparten el mismo sorteo del desenlace, el sesgo y la varianza de los
estimadores se incrementan y el tamaño de muestra efectivo se reduce
sustancialmente \cite{brill_winprob_2026}. Karbalaie \textit{et al.} midieron
el efecto sobre las cifras reportadas en un conjunto clínico de medidas
repetidas, donde un clasificador de árboles extremadamente aleatorizados
obtuvo 0.91 de acierto bajo validación cruzada estratificada de diez pliegues
y 0.66 bajo validación dejando un participante fuera, una inflación de 0.25
producida enteramente por la regla de partición
\cite{karbalaie_participantaware_2026}. Figueiredo y Mendes cuantifican el
mismo fenómeno en conjuntos de imágenes con alta correlación
espaciotemporal \cite{figueiredo_leakage_2024}.

El costo de ignorar la estructura no es, por tanto, un sesgo marginal: es la
diferencia entre un resultado y un artefacto.

\razon{Este párrafo debe existir en la introducción y no solo en la
metodología. El riesgo metodológico propio del \textit{track} seleccionado es
la fuga de información, y presentarlo desde el planteamiento del problema
convierte su control en parte de la contribución en lugar de en una
precaución técnica mencionada al final.}

\subsubsection{El problema es también de procedencia}

Puesto que el corpus fue recolectado por un tercero y publicado en medio de
una competencia en curso, no controlamos su procedencia, que es precisamente
la situación que Kaufman \textit{et al.} identifican como el caso más difícil
para el control de la fuga de datos \cite{kaufman_leakage_2011}. Establecer
que una cifra reportada significa lo que aparenta significar forma parte, por
tanto, del problema de investigación y no de sus preliminares.

El problema metodológico central consiste entonces en estimar
$P(\mathrm{victoria}\mid\mathrm{estado})$ empleando exclusivamente
información observable hasta el instante de predicción, y manteniendo
intactas las unidades independientes durante la evaluación.

\subsection{Pregunta de investigación}
\label{subsec:pregunta}

\begin{quote}
\textbf{PI:} ¿En qué medida las características observables del estado de una
partida de \emph{Pok\'emon TCG} permiten estimar la probabilidad de su
desenlace mediante modelos clásicos de aprendizaje automático; cómo varía esa
capacidad predictiva conforme progresa la partida; y cuánta se conserva al
evaluar sobre episodios, mazos y periodos temporales ausentes del
entrenamiento?
\end{quote}

\noindent\textbf{Condiciones de refutación.} La respuesta sería negativa, o
quedaría sustancialmente debilitada, si el intervalo de confianza del área
bajo la curva ROC en la ventana temprana incluyera el valor 0.5; si el
desempeño no mejorara al avanzar la partida una vez controlada su longitud;
si el modelo no superara a una heurística de dominio sobre partidas no
vistas; o si la señal desapareciera bajo el protocolo temporal o el de
reserva de mazos.

\noindent\textbf{Condiciones de medición.} La discriminación se mide mediante
el área bajo la curva ROC, y la calidad de la probabilidad mediante la
pérdida logarítmica, el puntaje de Brier y curvas de fiabilidad. Las métricas
se reportan de forma global y por ventana de turno. Las particiones y los
intervalos de confianza se construyen sobre episodios, aunque cada fila de la
representación tabular corresponda a una observación individual.

\noindent\textbf{Condiciones de viabilidad.} La pregunta puede responderse
dentro del semestre: los datos ya están descargados y auditados, y definir el
protocolo no requiere entrenar ningún modelo.

\razon{El enunciado exige que la pregunta pueda responderse con evidencia y
admita refutación. Enumerar las condiciones concretas bajo las cuales la
respuesta sería negativa es lo que distingue una pregunta de investigación de
una expectativa.}

\subsection{Hipótesis de trabajo}
\label{subsec:hipotesis}

\begin{quote}
\textbf{H1.} Las características estructuradas del estado observable
contienen señal predictiva sobre el desenlace desde etapas tempranas de la
partida, área bajo la curva ROC significativamente superior a 0.5 en la
ventana temprana, y esa señal aumenta conforme la partida progresa.
\end{quote}

\begin{quote}
\textbf{H2.} La señal temprana no se explica únicamente por el orden de turno
ni por la identidad del mazo.
\end{quote}

\begin{quote}
\textbf{H3.} Bajo el protocolo temporal, la calibración se degrada antes y en
mayor medida que la discriminación.
\end{quote}

\noindent\textbf{Fundamento de H3.} No es una conjetura sino una predicción
importada de un campo vecino. Guo \textit{et al.} revisaron quince estudios
clínicos sobre cambio temporal del conjunto de datos y encontraron que los
once estudios que midieron calibración reportaron deterioro de la
calibración, mientras que solo tres de los doce que midieron discriminación
reportaron deterioro de la discriminación \cite{guo_temporalshift_2021}.
Xenopoulos \textit{et al.} observan el mismo patrón en un juego competitivo:
un modelo evaluado fuera de su entorno presenta un incremento moderado en la
pérdida logarítmica acompañado de una degradación considerablemente mayor en
el error de calibración esperado \cite{xenopoulos_winprob_2022}. Silva Filho
\textit{et al.} ofrecen la explicación formal: las reglas de puntuación
propias se descomponen en un componente de calibración y uno de refinamiento,
y las métricas basadas exclusivamente en el ordenamiento son invariantes ante
transformaciones monótonas de las probabilidades predichas
\cite{silvafilho_calibration_2023}.

Que esa asimetría se sostenga en un dominio de juegos es una pregunta
empírica abierta, y enunciarla como hipótesis la vuelve falsable aquí.

\noindent\textbf{Criterios de falsación.} La componente temprana de H1 se
debilita si el intervalo de confianza del área bajo la curva en esa ventana
incluye 0.5, y su componente dinámica queda refutada si el valor por ventana
no crece una vez controlada la duración de la partida. H2 queda refutada si un
modelo restringido al orden de turno y a la identidad del mazo iguala al
conjunto completo de características, lo que indicaría un artefacto de
emparejamiento y no señal del estado de juego. H3 queda refutada si el puntaje
de Brier y las curvas de fiabilidad no se degradan más rápido que el área bajo
la curva bajo la partición temporal. Una caída hasta el nivel del azar en el
protocolo temporal se interpretará como ausencia de generalización y no como
fallo de ejecución.

No se fija de antemano ningún umbral numérico de mejora: sin resultado
empírico previo sobre este corpus, cualquier número de ese tipo sería
inventado.

\noindent\textbf{Advertencia sobre H1.} Brill \textit{et al.} documentan que
el sesgo del estimador es considerablemente mayor al comienzo de la partida,
por dos razones estructurales: los estados tempranos con desequilibrios
marcados son infrecuentes, y resulta más difícil vincular un estado temprano
con un desenlace determinado al final \cite{brill_winprob_2026}. En
consecuencia, un desempeño reducido en la ventana temprana admite dos
interpretaciones que este diseño debe procurar distinguir: ausencia genuina
de señal, o sesgo del estimador por escasez de observaciones en esa región
del espacio de estados.

\razon{Declarar esta ambigüedad antes de obtener resultados es obligatorio.
Sin ella, un desempeño bajo en turnos tempranos se interpretaría como
refutación de H1 cuando podría deberse al sesgo que la referencia documenta.}

\decidir{Para separar ambas interpretaciones hay que definir en la Etapa 2 un
procedimiento adicional: reportar el conteo de episodios y de observaciones
por ventana de turno junto a las métricas, o restringir el análisis temprano
a regiones del espacio de estados con soporte suficiente. La decisión afecta
el diseño experimental y conviene tomarla antes de ejecutar.}

\subsection{Objetivos}
\label{subsec:objetivos}

\noindent\textbf{Objetivo general.} Desarrollar y evaluar modelos clásicos de
aprendizaje automático que estimen la probabilidad de victoria en partidas de
\emph{Pok\'emon TCG} a partir de características tabulares derivadas
exclusivamente del estado observable, y caracterizar cómo evoluciona esa
capacidad predictiva a lo largo de la partida y ante cambios de episodio,
mazo y periodo temporal.

\noindent\textbf{Objetivos específicos.}
\begin{enumerate}
  \item Formalizar la estimación de la probabilidad de victoria a partir del
        estado de juego como un problema de clasificación binaria
        probabilística sobre datos tabulares, estableciendo la distinción
        entre la unidad de observación y la unidad estadísticamente
        independiente.
  \item Caracterizar el corpus y su calidad sobre la representación tabular
        final: análisis exploratorio, diccionario de datos, mecanismos de
        ausencia, valores atípicos, balance de clases y representatividad.
  \item Construir una tubería reproducible que vaya de la repetición en
        formato JSON a una tabla de instantáneas frescas con perspectiva
        normalizada, semillas fijas y verificación de esquema.
  \item Diseñar y verificar los controles de fuga de datos: agrupación por
        episodio, lista negra explícita de campos prohibidos y pruebas
        automatizadas que fallen si una columna prohibida entra al vector de
        características.
  \item Comparar una escalera de modelos de referencia, desde un predictor
        trivial hasta potenciación del gradiente, bajo tres protocolos de
        partición, reportando discriminación y calibración con intervalos de
        confianza calculados sobre episodios.
  \item Analizar explicabilidad y generalización: ablaciones por familia de
        características, importancia por permutación agrupada, curvas de
        fiabilidad y la curva de capacidad predictiva por turno.
\end{enumerate}

\razon{Los objetivos deben poder verificarse al final del semestre. Se
redactan como acciones con producto observable, formalizar, caracterizar,
construir, diseñar, comparar, analizar, y no como intenciones de desempeño,
porque un objetivo del tipo «superar un umbral» quedaría incumplido si el
resultado fuera negativo, cuando un resultado negativo también responde la
pregunta.}

\subsection{Contribución esperada}
\label{subsec:contribucion}

Se esperan tres contribuciones.

La primera es una caracterización de la probabilidad de victoria desde el
estado de juego en un juego de cartas que no ha sido estudiado de esta forma,
incluyendo el perfil por turno de la capacidad predictiva que Janusz
\textit{et al.} reportan para \emph{Hearthstone} pero que no existe para el
\emph{Pok\'emon TCG} \cite{janusz_hearthstone_2017}.

La segunda es un protocolo de evaluación cuyos controles de fuga están
medidos y no solo afirmados: la auditoría cuantifica la contaminación que
produciría el protocolo ingenuo, de modo que el diseño agrupado se justifica
con una cifra de este corpus y no por analogía.

La tercera es un análisis explícito de la diferencia entre discriminación y
calidad probabilística, dimensión que Grinsztajn \textit{et al.} identifican
como pendiente en la evaluación de modelos tabulares
\cite{grinsztajn_trees_2022} y que la literatura de predicción de desenlace
en juegos ha reportado de forma desigual. A ello se suma una tubería
reproducible y un diccionario de datos documentado para un corpus que hoy
carece de descripción arbitrada, siguiendo la estructura que el descriptor
SC2EGSet establece para datos de repeticiones de deportes electrónicos
\cite{bialecki_sc2egset_2023}.

Un resultado negativo también sería una contribución. Si la discriminación
colapsa bajo el protocolo temporal, eso constituye evidencia sobre la rapidez
con que un conjunto de estrategias dominante invalida un evaluador de
estados, que es la pregunta que Bertram \textit{et al.} plantean para
repertorios de cartas en cambio permanente
\cite{bertram_cardrepresentations_2024}.

\razon{La contribución se enuncia en términos de aportes verificables y no de
resultados prometidos. Comprometerse en la Etapa 1 con un nivel de desempeño
ata el trabajo a un resultado que aún no se conoce.}

\subsection{Estructura del documento}
\label{subsec:estructura}

La Sección~\ref{sec:relacionados} revisa la literatura en tres ejes, el
problema y su dominio, las técnicas del \textit{track} seleccionado, y el
conjunto de datos, y cierra con la identificación explícita de las brechas.
La Sección~\ref{sec:baseline} declara el punto de referencia tomado de la
literatura y discute su comparabilidad. La Sección~\ref{sec:metodologia}
presenta la formalización del problema, el protocolo experimental y el plan
de modelos de referencia. La Sección~\ref{sec:analisistecnico} desarrolla el
análisis técnico exigido por el \textit{track} seleccionado, con énfasis en
la prevención de fuga de información. La Sección~\ref{sec:etica} documenta
las consideraciones éticas preliminares.